\begin{figure}[!htbp]
\centering
\begin{tikzpicture}[x=1cm,y=1cm]
\node[paneltitle] at (-2,2.25) {(a) Partition into three arcs};
\node[paneltitle] at (6,2.25) {(b) Multiply within each arc};
\foreach \i/\a in {0/120,1/60,2/0,3/-60,4/-120,5/180}{\node[observer,minimum size=7mm] (o\i) at (\a:1.5) {$O_\i$};}
\begin{scope}[on background layer]
\draw[inkblue!15,line width=15pt] (o0)--(o1);
\draw[inkrust!15,line width=15pt] (o2)--(o3);
\draw[inkgreen!15,line width=15pt] (o4)--(o5);
\draw[network,draw=inkblue,line width=1.3pt] (o0)--(o1);
\draw[network,draw=inkrust,line width=1.3pt] (o2)--(o3);
\draw[network,draw=inkgreen,line width=1.3pt] (o4)--(o5);
\draw[network] (o1)--(o2) (o3)--(o4) (o5)--(o0);
\end{scope}
\node[fignote] at (1.48,.88) {$S_1$};
\node[fignote] at (0,-1.72) {$S_3$};
\node[fignote] at (-1.6,.85) {$S_5$};
\draw[-{Latex[length=2mm]},black!60] (2.45,0)--(5.3,0);
\node[fignote] at (3.85,.6) {Boundary sources\\remain independent};
\node[fignote] at (3.85,-.65) {Internal sources become\\private randomness};
\node[observer,draw=inkblue,fill=inkblue!8] (v1) at (8,1.25) {$V_1$};
\node[observer,draw=inkrust,fill=inkrust!8] (v2) at (9.6,-1.25) {$V_2$};
\node[observer,draw=inkgreen,fill=inkgreen!8] (v3) at (6.4,-1.25) {$V_3$};
\draw[network] (v1)--node[right=2pt,font=\small] {$S_1$}(v2)--node[below=2pt,font=\small] {$S_3$}(v3)--node[left=2pt,font=\small] {$S_5$}(v1);
\node[fignote,inkblue] at (8,1.8) {$O_0O_1$};
\node[fignote,inkrust] at (10.5,-1.25) {$O_2O_3$};
\node[fignote,inkgreen] at (5.6,-1.65) {$O_4O_5$};
\node[figlabel] at (0,-2.5) {Example: $C_6$};
\node[figlabel] at (8,-2.5) {$V_r=\prod_{v\in\Lambda_r}O_v$};
\end{tikzpicture}
\caption{Reducing cycle incompatibility to triangle parity rigidity. Partition any cycle into three nonempty contiguous arcs and multiply the observations within each arc. The three boundary sources form a triangle; internal sources become private randomness. The target reduces to $\Pi(-q,-q,-q)$, and total variation cannot increase under this map.}\label{fig:cycle-reduction}
\end{figure}
